\chapter{CODE ATTACHMENTS}
\label{app:code}

This appendix reproduces short excerpts of the project source that implement the parts of the module evaluated in this report. The lines are copied verbatim from the repository; the numbers in the margin are the line numbers in the original file. Excerpts are limited to the code discussed in the chapters, so the rest of each file is omitted.

\section{Relationship label schema}
\label{app:A1}
The eight labels and the definitions given to the model (Section~2.2, Table~2.1). The reviewer form uses the same eight labels. Source: \code{backend/\allowbreak{}app/\allowbreak{}knowledge\_engine/\allowbreak{}semantic\_provider.py}, lines 33--44.

\begin{lstlisting}[style=repo, firstnumber=33]
RELATIONSHIPS = ("SAME_CONCEPT", "PARENT", "CHILD", "RELATED", "PARTIAL_OVERLAP", "UNRELATED", "AMBIGUOUS",
                 "INSUFFICIENT_EVIDENCE")
RELATIONSHIP_DEFINITIONS = {
    "SAME_CONCEPT": "LEFT and RIGHT denote the same thing in this context",
    "PARENT": "LEFT is broader than RIGHT (RIGHT is a kind or part of LEFT)",
    "CHILD": "LEFT is narrower than RIGHT (LEFT is a kind or part of RIGHT)",
    "RELATED": "clearly associated, neither contains the other",
    "PARTIAL_OVERLAP": "they share part of their meaning",
    "UNRELATED": "different meanings (for example a different word sense)",
    "AMBIGUOUS": "the context supports more than one reading",
    "INSUFFICIENT_EVIDENCE": "the text does not allow a judgement",
}
\end{lstlisting}

\section{Ordered decision procedure of the v4 prompt experiment}
\label{app:A2}
The decision procedure added in prompt adjudicate-v4 (Section~4.1.2). It was evaluated and not adopted; production uses adjudicate-v3. Source: \code{backend/\allowbreak{}app/\allowbreak{}knowledge\_engine/\allowbreak{}semantic\_v3.py}, lines 203--213.

\begin{lstlisting}[style=repo, firstnumber=203]
DECISION_PROCEDURE_V4 = (
    "Decide in this order:\n"
    "1. Does LEFT, as used in the SOURCE CONTEXT, mean the RIGHT concept or something connected to it? If the source "
    "uses LEFT in another sense, answer UNRELATED. If the source does not let you tell, answer AMBIGUOUS (several "
    "readings fit) or INSUFFICIENT_EVIDENCE (nothing decides it).\n"
    "2. SAME_CONCEPT only if LEFT and RIGHT are interchangeable here: the same thing at the same breadth.\n"
    "3. CHILD if LEFT is one specific kind, instance, part, product, platform, vendor or tool of RIGHT. A named "
    "product or provider of a technology category is CHILD of that category, not SAME_CONCEPT.\n"
    "4. PARENT if LEFT is a broader area or category that includes RIGHT as one of its kinds or parts.\n"
    "5. PARTIAL_OVERLAP if they share part of their meaning but each also covers things the other does not.\n"
    "6. RELATED if they are associated (used together, a tool for, a prerequisite of) without sharing meaning.\n")
\end{lstlisting}

\section{Deterministic evidence validator}
\label{app:A3}
The rules every model answer must pass before it can reach the analysis (Section~4.1.1). Lines 31 to 279 of the file are omitted. Source: \code{backend/\allowbreak{}app/\allowbreak{}knowledge\_engine/\allowbreak{}semantic\_v3.py}, lines 28--30, 280--354.

\begin{lstlisting}[style=repo, firstnumber=28]
METADATA_TYPES = ("TAXONOMY_METADATA", "NORMALIZATION_METADATA")
DEFINITE = ("SAME_CONCEPT", "PARENT", "CHILD", "RELATED", "PARTIAL_OVERLAP", "UNRELATED")
MEANING_RELATIONS = ("SAME_CONCEPT", "PARENT", "CHILD", "RELATED", "PARTIAL_OVERLAP")
\end{lstlisting}

\noindent{\small\itshape Lines 31--279 omitted.}

\begin{lstlisting}[style=repo, firstnumber=280]
DIRECTIONAL_TYPES = ("ESCO_HIERARCHY", "ESCO_RELATION", "COURSE_TITLE", "UNIT_TITLE", "LEARNING_OUTCOME")


def validate_decision(out, c):
    """
    Deterministic gate over a model decision. Returns (final_state, findings).

    final_state:
      MODEL_SUPPORTED        definite relationship whose evidence passes every rule
      MODEL_SUGGESTED        AMBIGUOUS / INSUFFICIENT_EVIDENCE answered by the model (kept for review, no effect)
      REVIEW_REQUIRED        schema failure or evidence that cannot be traced to the inputs
      INSUFFICIENT_EVIDENCE  definite relationship without the required source evidence (downgraded)
    Rules: evidence present; each cited ref exists; the quote occurs in that item; the source type is OUR item's type
    (a model-claimed type that differs is recorded); metadata items never count; any definite relationship needs a
    SOURCE-side item; meaning relations (SAME/PARENT/CHILD/RELATED/PARTIAL) also need a CONCEPT-side item;
    PARENT/CHILD need a direction justification AND a cited directional item (DIRECTIONAL_TYPES); SAME_CONCEPT on alias/label candidates needs SOURCE_TEXT or JOB_TITLE.
    """
    findings = []
    schema = validate_output(out)
    if schema:
        return "REVIEW_REQUIRED", [f"SCHEMA: {e}" for e in schema]
    rel = out["relationship"]
    if rel not in DEFINITE:
        return "MODEL_SUGGESTED", ["NON_DEFINITE_RELATIONSHIP"]
    items = {it["id"]: ("LEFT", it) for it in c["left_items"]}
    items.update({it["id"]: ("RIGHT", it) for it in c["right_items"]})
    items.update({it["id"]: ("META", it) for it in c["metadata_items"]})
    valid = []
    for e in out.get("evidence", []):
        side_it = items.get(e["ref"])
        if side_it is None:
            findings.append(f"UNKNOWN_REF:{e['ref']}")
            continue
        side, it = side_it
        if len(_norm(e["quote"])) < 2:
            findings.append(f"EMPTY_QUOTE:{e['ref']}")
            continue
        if _norm(e["quote"]) not in _norm(it["text"]):
            findings.append(f"QUOTE_NOT_IN_ITEM:{e['ref']}")
            continue
        if e.get("source_type") and e.get("source_type") != it["source_type"]:
            findings.append(f"SOURCE_TYPE_MISLABELLED:{e['ref']}:{e.get('source_type')}->{it['source_type']}")
        if side == "META" or it["source_type"] in METADATA_TYPES:
            findings.append(f"METADATA_CITED:{e['ref']}")
            continue
        valid.append((side, it["source_type"]))
    untraceable = any(f.startswith(("UNKNOWN_REF", "QUOTE_NOT_IN_ITEM")) for f in findings)  # invented evidence
    has_left = any(s == "LEFT" for s, _ in valid)
    has_right = any(s == "RIGHT" for s, _ in valid)
    if not out.get("evidence"):
        findings.append("NO_EVIDENCE")
    if not has_left:
        findings.append("NO_SOURCE_EVIDENCE")
    if rel in MEANING_RELATIONS and not has_right:
        findings.append("NO_CONCEPT_EVIDENCE")
    if rel == "SAME_CONCEPT" and c["kind"] in ("ALIAS_SENSE", "AMBIGUOUS_TAXONOMY_LABEL") \
            and not any(s == "LEFT" and t in ("SOURCE_TEXT", "JOB_TITLE") for s, t in valid):
        findings.append("ALIAS_SAME_CONCEPT_WITHOUT_SOURCE_TEXT")
    if rel in ("PARENT", "CHILD") and not (out.get("direction_justification") or "").strip():
        findings.append("NO_DIRECTION_JUSTIFICATION")
    if rel in ("PARENT", "CHILD") and not any(t in DIRECTIONAL_TYPES for _, t in valid):
        # a shared word or one string inside another is never hierarchy: direction needs a structural fact
        findings.append("NO_DIRECTIONAL_EVIDENCE")
    blocking = {"NO_EVIDENCE", "NO_SOURCE_EVIDENCE", "NO_CONCEPT_EVIDENCE", "ALIAS_SAME_CONCEPT_WITHOUT_SOURCE_TEXT",
                "NO_DIRECTION_JUSTIFICATION", "NO_DIRECTIONAL_EVIDENCE"}
    if untraceable:
        return "REVIEW_REQUIRED", findings
    if any(f in blocking for f in findings):
        return "INSUFFICIENT_EVIDENCE", findings
    return "MODEL_SUPPORTED", findings


def effective_relationship(out, final_state):
    """What downstream may read: the model's relationship only when MODEL_SUPPORTED."""
    return out["relationship"] if final_state == "MODEL_SUPPORTED" else final_state
\end{lstlisting}

\section{Reading a model decision as downstream would}
\label{app:A4}
How a recorded decision becomes the validated answer that is scored in Chapters 4 and 5 (Section~3.2.2). Source: \code{scripts/\allowbreak{}semantic/\allowbreak{}score\_against\_reviewers.py}, lines 42--53.

\begin{lstlisting}[style=repo, firstnumber=42]
def validated(d, cand):
    """What downstream may read (same rule as scripts/semantic/evaluate_v3.py `effective`)."""
    if d is None:
        return None, ["NO_DECISION"]
    if d.get("decision_state") not in ("MODEL_SUGGESTED", "MODEL_SUPPORTED", "INSUFFICIENT_EVIDENCE", "REVIEW_REQUIRED"):
        return "AMBIGUOUS", [f"DECISION_STATE:{d.get('decision_state')}"]
    out = {k: d.get(k) for k in ("relationship", "reason", "evidence", "direction_justification", "alternatives",
                                 "uncertainty")}
    state, findings = v3.validate_decision(out, cand)
    if state in ("MODEL_SUPPORTED", "MODEL_SUGGESTED"):
        return d["relationship"], findings
    return ("INSUFFICIENT_EVIDENCE" if state == "INSUFFICIENT_EVIDENCE" else "AMBIGUOUS"), findings
\end{lstlisting}

\section{Cohen's kappa between two reviewers}
\label{app:A5}
The agreement statistic of Equation~3.4, used for Table~5.6. Source: \code{backend/\allowbreak{}app/\allowbreak{}knowledge\_engine/\allowbreak{}semantic\_eval.py}, lines 47--54.

\begin{lstlisting}[style=repo, firstnumber=47]
def cohen_kappa(a, b):
    n = len(a)
    if not n:
        return None
    po = sum(x == y for x, y in zip(a, b)) / n
    ca, cb = Counter(a), Counter(b)
    pe = sum(ca[k] * cb[k] for k in set(ca) | set(cb)) / (n * n)
    return None if pe == 1 else round((po - pe) / (1 - pe), 4)
\end{lstlisting}

\section{Applying the promotion gate}
\label{app:A6}
Every criterion of semantic-gate-v3 (Table~3.3) is checked; any failure leaves the system in SAFE\_MODE (Section~3.2.5). Source: \code{backend/\allowbreak{}app/\allowbreak{}knowledge\_engine/\allowbreak{}semantic\_eval.py}, lines 257--292.

\begin{lstlisting}[style=repo, firstnumber=257]
def apply_gate(gate, gold_rows, reviewers_agreement, system_result, label_sets):
    """Every check with its value and threshold; promotable classes only when the whole gate passes."""
    gold_cases = sum(1 for r in gold_rows if r["gold"])
    c = gate["criteria"]
    checks = {}

    def check(name, value, ok, threshold):
        checks[name] = {"value": value, "threshold": threshold, "pass": bool(ok)}
    check("independent_reviewers", len(label_sets), len(label_sets) >= c["min_independent_reviewers"], c["min_independent_reviewers"])
    check("gold_cases", gold_cases, gold_cases >= c["min_gold_cases"], c["min_gold_cases"])
    k = reviewers_agreement.get("kappa")
    if len(label_sets) >= 2:
        check("inter_rater_kappa", k, k is not None and k >= c["min_kappa_when_two_or_more_reviewers"],
              c["min_kappa_when_two_or_more_reviewers"])
    s = system_result
    check("accuracy", s["accuracy"], (s["accuracy"] or 0) >= c["min_accuracy"], c["min_accuracy"])
    check("macro_f1", s["macro_f1"], (s["macro_f1"] or 0) >= c["min_macro_f1"], c["min_macro_f1"])
    e = s["critical_errors"]
    check("false_same_concept", e.get("FALSE_SAME_CONCEPT", 0), e.get("FALSE_SAME_CONCEPT", 0) <= c["max_false_same_concept"], c["max_false_same_concept"])
    check("wrong_hierarchy_direction", e.get("WRONG_HIERARCHY_DIRECTION", 0),
          e.get("WRONG_HIERARCHY_DIRECTION", 0) <= c["max_wrong_hierarchy_direction"], c["max_wrong_hierarchy_direction"])
    fh = e.get("FALSE_PARENT", 0) + e.get("FALSE_CHILD", 0)
    check("false_parent_or_child", fh, fh <= c["max_false_parent_or_child"], c["max_false_parent_or_child"])
    check("unsupported_taxonomy_evidence", e.get("UNSUPPORTED_TAXONOMY_EVIDENCE", 0),
          e.get("UNSUPPORTED_TAXONOMY_EVIDENCE", 0) <= c["max_unsupported_taxonomy_evidence"], c["max_unsupported_taxonomy_evidence"])
    check("lexical_collision", e.get("LEXICAL_COLLISION", 0), e.get("LEXICAL_COLLISION", 0) <= c["max_lexical_collision"], c["max_lexical_collision"])
    passed = all(v["pass"] for v in checks.values())
    promotable = []
    if passed:
        for cls in gate["promotable_classes"]:
            pc = s["per_class"][cls]
            if pc["precision"] is not None and pc["predicted"] >= c["min_predictions_per_promoted_class"] \
                    and pc["precision"] >= c["min_precision_per_promoted_class"]:
                promotable.append(cls)
    return {"passed": passed, "checks": checks, "promotable_classes": promotable,
            "mode": "PROMOTE_VALIDATED_CLASSES" if passed and promotable else "SAFE_MODE"}
\end{lstlisting}

\section{Storing an independent label set only with an attestation}
\label{app:A7}
An independent label set cannot be created without a statement that the reviewer answered alone and without AI tools (Section~4.3.4). Source: \code{scripts/\allowbreak{}evaluation/\allowbreak{}semantic\_validation.py}, lines 709--731.

\begin{lstlisting}[style=repo, firstnumber=709]
def _store_labels(pkg, man, reviewer, rows, source_file, meta=None, attestation=None):
    """An INDEPENDENT_HUMAN_REVIEWER label set: it counts for agreement, gold and the gate, so it is stored only with
    a statement that the reviewer answered alone, from the case text, without AI tools and without seeing system
    output or another reviewer's answers (attestation: who confirmed it, and how). Responses that do not meet this go
    through import-responses instead."""
    try:
        name = se.check_reviewer(reviewer)
    except ValueError as exc:
        raise SystemExit(str(exc)) from exc
    if not (attestation or "").strip():
        raise SystemExit("an independent label set needs --attestation: who confirmed that the reviewer answered alone, "
                         "without AI tools, and how; otherwise import the responses with import-responses")
    out = pkg / "labels" / f"labels_{slug(name)}_v1.json"
    if out.exists():
        raise SystemExit(f"{out.name} exists; label sets are immutable")
    doc = {"benchmark": man["benchmark_version"], "benchmark_sha256": man["benchmark"]["sha256"], "package": man["package"],
           "reviewer": name, "reviewer_background": meta or {}, "label_version": "v1", "label_schema_version": man["label_schema_version"],
           "independent": True, "model_blind": True, "label_source": "INDEPENDENT_HUMAN_REVIEWER",
           "attestation": {"statement": attestation.strip(), "recorded_at": now()},
           "ingested_at": now(), "source_file_sha256": sha(Path(source_file).read_bytes()),
           "labels": rows, "labelled": sum(1 for r in rows if r["label"])}
    out.write_text(json.dumps(doc, indent=1, ensure_ascii=False))
    return out, doc
\end{lstlisting}

